\section*{\texorpdfstring{\S\CurrentExerciseGroup}{第{\CurrentExerciseGroup}节}}
\phantomsection
\addcontentsline{toc}{section}{第{\CurrentExerciseGroup}节习题}

\begin{enumerate}
\def\labelenumi{\arabic{enumi})}
\tightlist
\item
  设 E 是一个 Riesz 空间，x 是 E 的一个 \textgreater0 的元素。若存在 E 上的一个相对有界线性形式 L，使得 \(L(x) \neq 0\) ，试证 nx（n 为整数 \textgreater0）所成的集在 E 中不可能是上有界的。特别地，试证在 §1 习题 5 b) 所定义的 Riesz 空间 E 上不存在任何相对有界线性形式。
\end{enumerate}

¶ 2) a) 设 L 是 Riesz 空间 E 上的一个正线性形式。考虑 E 的一个元素 a \textgreater{} 0 以及满足下列条件的正线性形式 \(M \leqslant L\) 所成的集：\(1^{\circ} M(x) = L(x)\) 对每个满足 \(0 \leqslant x \leqslant a\) 的 x 成立；\(2^{\circ} M(x) = 0\) 对每个与 a 互奇异的 \(x \geqslant 0\) 成立。试证这个正线性形式的集有一个最大元素 \(L_{a}\) ，且对每个 \(x \geqslant 0\) ，\(L_{a}(x) = \inf L(y)\) ，其中 y 取遍满足 \(0 \leqslant y \leqslant x\) 且 x - y 与 a 互奇异的一切元素所成的集。试证 \(L_{\lambda a} = L_{a}\) 对每个标量 \(\lambda > 0\) 成立，且若 a 与 b 互奇异，则 \(L_{a+b} \leqslant L_{a} + L_{b}\) 。若 E 是完全格序的，试证当 a 与 b 在 E 中互奇异时 \(L_{a+b} = L_{a} + L_{b}\) 。

\begin{enumerate}
\def\labelenumi{\alph{enumi})}
\setcounter{enumi}{1}
\tightlist
\item
  取 E 为区间 \(\mathbf{I} = [0,1]\) （属于 \(\mathbf{R}\) ）上连续实值函数所成的 Riesz 空间，并取 \(L(x) = x\left(\frac{1}{2}\right)\) ；试用一个例子证明可以有 \(L_{a + b} < L_a + L_b\) ，对两个互奇异元素 \(a, b\) （属于 \(\mathbf{E}\) ）而言。
\end{enumerate}

\begin{enumerate}
\def\labelenumi{\arabic{enumi})}
\setcounter{enumi}{2}
\tightlist
\item
  设 \(\mathbf{E}\) 是一个完全格序空间。
\end{enumerate}

\begin{enumerate}
\def\labelenumi{\alph{enumi})}
\item
  试证：若线性形式 \(L\) 在 \(\mathbf{E}\) 上对拓扑 \(\mathcal{T}_0(\mathbf{E})\) （§1, 习题 10）连续，则它相对有界，且 \(|L|\) 对 \(\mathcal{T}_0(\mathbf{E})\) 连续（用反证法，注意对每个 \(x \in \mathbf{E}\) ，元素序列 \(x / n\) 当 \(n\) 趋于无穷时对拓扑 \(\mathcal{T}_0(\mathbf{E})\) 趋于 0）。
\item
  设 A 是任一无穷集，U 是 A 上的一个超滤子，其诸集合的交为空。在完全格序空间 \(E = \mathcal{B}(A)\) 上，考虑正线性形式 \(x \mapsto \lim_{\mu} x(t)\) ；试证此线性形式对拓扑 \(\mathcal{T}_{0}(E)\) 不连续。
\end{enumerate}

\begin{enumerate}
\def\labelenumi{\arabic{enumi})}
\setcounter{enumi}{3}
\tightlist
\item
  设 E 是一个完全格序空间。
\end{enumerate}

\begin{enumerate}
\def\labelenumi{\alph{enumi})}
\item
  设 \(L\) 是 \(\mathbf{E}\) 上的一个正线性形式，对拓扑 \(\mathcal{T}_0(\mathbf{E})\) （§1, 习题 10）连续。试证由 \(x \in \mathbf{E}\) （满足 \(L(|x|) = 0\) ）组成的集是一个带 \(Z(L)\) 。设 \(S(L)\) 为与 \(Z(L)\) 在 \(\mathbf{E}\) 中互余的带（§1, 习题 6）。
\item
  设 L 与 M 是 E 上两个正线性形式，对 \(\mathcal{T}_{0}(\mathrm{E})\) 连续。为使 M 属于 E 上相对有界线性形式所成的空间 \(\Omega\) 中由 L 生成的带，必须且只须 \(S(M) \subset S(L)\) 。（为看出条件是充分的，用反证法，假设 M 不满足命题 5 的条件；考虑 E 的元素的一个序列 \((y_{n})\) ，满足 \(0 \leqslant y_{n} \leqslant x\) ，
\end{enumerate}

\(L(y_{n}) \leqslant 1 / 2^{n}\) 且 \(M(y_{n}) \geqslant \alpha > 0\) ，并由此推出存在 E 的一个元素 \(z \geqslant 0\) ，使得 \(L(z) = 0\) 且 \(M(z) \geqslant \alpha\) 。

\begin{enumerate}
\def\labelenumi{\alph{enumi})}
\setcounter{enumi}{2}
\tightlist
\item
  设 \(L\) 与 \(M\) 是 \(\mathbf{E}\) 上两个正线性形式，对 \(\mathcal{T}_0(\mathbf{E})\) 连续。为使在 \(\Omega\) 中 \(L\) 与 \(M\) 互奇异，必须且只须 \(S(L) \cap S(M) = \{0\}\) 。（为看出条件是必要的，试证若 \(S(L) \cap S(M)\) 不化为 0，则 \(\Omega\) 中由 \(L\) 与 \(M\) 生成的带有公共元素 \(\neq 0\) ，办法是利用 b）。）
\end{enumerate}

¶ 5) a) 设 E 是一个 Riesz 空间，H 是 E 的一个孤立线性子空间（§1, 习题 4）。设 Ω 是 E 上相对有界线性形式所成的空间，设 Θ 是 Ω 中由在 H 上为零的线性形式 \(L \in \Omega\) 组成的子空间。试证 Θ 是 Ω 的一个孤立子空间，且 Θ 是一个完全格序空间，同构于 Riesz 空间 E/H 上相对有界线性形式所成的空间。

\begin{enumerate}
\def\labelenumi{\alph{enumi})}
\setcounter{enumi}{1}
\tightlist
\item
  设 \(u\) 是 \(E\) 到一个 Riesz 空间 \(F\) 中的线性映射；下列条件等价：
\end{enumerate}

\(\alpha) u(\sup(x, y)) = \sup(u(x), u(y));\)

\(\beta) u(\inf(x, y)) = \inf(u(x), u(y));\)

\(\gamma) u(x^{+}) = (u(x))^{+}\) ;

\(\delta)\) 关系 \(\inf (x,y) = 0\) 蕴涵 \(\inf (u(x),u(y)) = 0\)

这时称 u 是一个格线性映射。

\begin{enumerate}
\def\labelenumi{\alph{enumi})}
\setcounter{enumi}{2}
\item
  为使 E 的一个线性子空间 H 在孤立子空间 \(\neq\) E 的集合中是极大的，必须且只须它形如 \(\operatorname{Ker}(f)\) ，其中 \(f\) 是一个格线性形式 \(\neq 0\) 。（利用 GT, V, §3, 习题 1。）
\item
  给出一个不化为 0 且不含任何极大孤立子空间的 Riesz 空间的例子（参见 §1, 习题 5 b)）。
\end{enumerate}

¶ 6) 设 E 是一个 Riesz 空间，Ω 是 E 上相对有界线性形式所成的完全格序空间，F 是 Ω 的一个余格子子空间（§1, 习题 3）。对每个 \(x \in E\) ，映射 \(x' \mapsto \langle x, x' \rangle\) （由 F 到 R 中）是 F 上的一个相对有界线性形式 \(u_x\) ，且映射 \(x \mapsto u_x\) 是 E 到 F 上相对有界线性形式所成的完全格序空间 Ω' 中的一个递增线性映射。为使 \(x \mapsto u_x\) 是 E 到 Ω' 的一个余格子子空间上的同构，也就是说，为使 \(u_x > 0\) 蕴涵 x \textgreater{} 0，且 \(u_{\text{sup}(x,y)} = \text{sup}(u_x, u_y)\) ，必须且只须下列两个条件成立：1° 对 E 中每个 x \textgreater{} 0，存在 F 中的 \(x' > 0\) 使 \(\langle x, x' \rangle > 0\) ；2° 对 E 的每对互奇异元素 \(y \geqslant 0\) ， \(z \geqslant 0\) ，对每个数 \(\varepsilon > 0\) 以及 F 中每个 \(x' \geqslant 0\) ，存在 F 的两个元素 \(y' \geqslant 0\) ， \(z' \geqslant 0\) ，使得 \(x' = y' + z'\) 且 \(\langle y, y' \rangle + \langle z, z' \rangle \leqslant \varepsilon\) 。（注意：若第二个条件成立且 v 是 F 上的一个正线性形式，满足 \(v \geqslant u_x\) ，则有 \(v(x') \geqslant \langle x^+, x' \rangle - \varepsilon\) ，对 F 中每个 \(x' \geqslant 0\) 及每个 \(\varepsilon > 0\) 。）

若条件 \(2^{\circ}\) 成立而条件 \(1^{\circ}\) 不成立，则由元素 \(x \in E\) （满足 \(\langle x, x' \rangle = 0\) ，对一切 \(x' \in F\) ）组成的集是孤立线性子空间 \(H\) （在 \(E\) 中）。通过过渡到商，映射 \(x \mapsto u_x\) 这时定义 Riesz 空间 \(E / H\) 到 \(\Omega'\) 的一个余格子子空间上的一个同构。

试证当 \(\mathbf{F} = \Omega\) 时，上述条件 \(2^{\circ}\) 总成立（利用习题 2）。

\begin{enumerate}
\def\labelenumi{\arabic{enumi})}
\setcounter{enumi}{6}
\tightlist
\item
  设 \(E\) 是一个维数 \(n\) 有限的 Riesz 空间。
\end{enumerate}

\begin{enumerate}
\def\labelenumi{\alph{enumi})}
\item
  若 E 是阿基米德的（§1, 习题 5），试证 E 的元素 \(\geqslant 0\) 所成的锥 P 是闭的且有内点。由此推出 P 的诸支撑超平面的交化为 0（利用 \(P \cap (-P) = \{0\}\) 这一事实）。
\item
  利用 a) 与习题 6，试证每个维数 n 的阿基米德 Riesz 空间都同构于乘积空间 \(R^{n}\) （参见 §1, 习题 13 e)）。
\item
  给出一个维数 2 的全序 Riesz 空间的例子。
\end{enumerate}

¶ 8) 设 E 是一个 Riesz 空间，\((U_{\iota})\) 是 E 上正线性形式的一个族。我们考虑 E 上的拓扑 \(\mathcal{T}\) ，它由半范数 \(U_{\iota}(|x|)\) 定义。

\begin{enumerate}
\def\labelenumi{\alph{enumi})}
\item
  试证 E 中由满足 \(U_{\iota}(|x|)=0\) （对一切 \(\iota\) ）的 x 组成的子空间 H 是 E 的一个孤立子空间（§1, 习题 4）；E 相伴的 Hausdorff 空间是商空间 E/H；通过过渡到商，\(U_{\iota}\) 在 E/H 上定义正线性形式 \(\dot{U}_{\iota}\) ，且 E/H 上的商拓扑由半范数 \(\dot{U}_{\iota}(|\dot{x}|)\) 定义。
\item
  试证在空间 E 中映射 \(x \mapsto |x|\) 是一致连续的，并由此推出：若拓扑 T 是 Hausdorff 的，则它与 E 的序向量空间结构相容。
\item
  假设拓扑 \(\mathcal{T}\) 是 Hausdorff 的且 E 是完全格序的；试证 E 中每个带 B 都是闭的，且若 B' 是由与 B 的每个元素都互奇异的元素组成的带，则 E 是 B 与 B' 的拓扑直和。
\item
  假设拓扑 \(\mathcal{T}\) 是 Hausdorff 的。设 \(\overline{\mathrm{E}}\) 是 \(\mathrm{E}\) 的完备化；若 \(\mathrm{P}\) 是元素 \(\geqslant 0\) （属于 \(\mathrm{E}\) ）所成的集，试证 \(\overline{\mathrm{P}}\) （即 \(\mathrm{P}\) 在 \(\widehat{\mathrm{E}}\) 中的闭包）在 \(\widehat{\mathrm{E}}\) 上定义一个 Riesz 空间结构（参见 §1, No. 2）。
\item
  若 E 对拓扑 T 是 Hausdorff 的且完备的，试证为使一个集 \(A \subset E\) （按关系 \(\leqslant\) 有向）在 E 中有一个上确界，必须且只须对每个指标 \(\iota\) ，\(U_{\iota}(x)\) 在 A 中上有界；这时对 E 上每个连续的递增数值函数 f，有 \(\sup_{x \in A} f(x) = f(\sup A)\) 。由此
\end{enumerate}

推出 E 是完全格序的。

\begin{enumerate}
\def\labelenumi{\alph{enumi})}
\setcounter{enumi}{5}
\tightlist
\item
  假设 E 对拓扑 T 是 Hausdorff 的且完备的。试证若 E 上的滤子 F 关于 E 的序结构有一个极限（§1, 习题 9），则它对拓扑 T 收敛于同一极限（利用 e)）。
\end{enumerate}

\begin{enumerate}
\def\labelenumi{\arabic{enumi})}
\setcounter{enumi}{8}
\tightlist
\item
  设 E 是一个 Riesz 空间，\(\Omega\) 是 E 上相对有界线性形式所成的空间。考虑 \(\Omega\) 上由半范数 \(L \mapsto |L|(x)\) 定义的拓扑，其中 x 取遍 E 的元素 \(\geqslant 0\) 所成的集。试证赋以此拓扑的 \(\Omega\) 是 Hausdorff 的且完备的。
\end{enumerate}

